\chapter*{Abstract}

\noindent
A retrieval-augmented generation system can answer only from the passages its
retriever finds. On earnings call transcripts the same companies discuss the same
topics in similar language every quarter, so the hard part is finding the right
passage rather than a near-identical wrong one. Three retrieval-design choices were crossed over a fixed corpus,
question set, prompt, and generator: chunk size, retrieval strategy, and
indexing representation, meaning whether each passage is indexed as raw text
alone or with the company and quarter it comes from. Retrieval
and answers were scored separately. Chunk size made no detectable difference,
and retrieval strategy made none until metadata was indexed, while indexing
representation moved retrieval substantially. The first two had nothing to
work with: what tells these documents apart is the company and the quarter,
which never reached the text the retriever searched. Better retrieval did not
carry through: no difference in answer outcomes between conditions was
detectable, a failure to detect one at this sample size rather than evidence
that the conditions answer alike. Every answer the
automatic grounding check called ungrounded was read, and some of it was
faithful text the check could not recognize: such a check errs in a predictable
direction. When the retrieved context cannot support an answer,
abstaining is correct, and an evaluation should reward it. Effort spent on
how a collection is cut up and searched can miss what distinguishes its
documents, which already sits in their metadata.

\vspace{1em}
\noindent
\textit{\textbf{Keywords: }retrieval-augmented generation; indexing
representation; chunk size; earnings call transcripts; abstention}
